\chapter{Auswertung}\label{ch:auswertung}

% Hinweis: Konkrete Mess-Resultate werden nach den ersten messung_driver-
% Laeufen eingefuegt. Aktuell beschreibt dieses Kapitel die Auswertungs-
% Methodik und die Hypothesen-Definitionen.

\section{Auswertungs-Pipeline}

Pro Messreihen-Lauf erzeugt der \texttt{messung\_driver} unter
\texttt{Code/\_runs/<date>/<spec\_id>/} folgende Dateien:

\begin{itemize}
    \item \texttt{measurements.csv} --- pro Permutation eine Zeile mit
          Cycles, L1/L2/L3-Misses, Branches, Throughput
    \item \texttt{measurements.json} --- gleiche Daten als JSON
    \item \texttt{raw.bin} (optional) --- binary Records mit
          \texttt{kBinaryMagic 0xC0FFEE02}
\end{itemize}

Die Auswertung nutzt drei Tools aus \texttt{Diplomarbeit/Code/}:

\begin{enumerate}
    \item \texttt{binary\_to\_csv} --- deserialisiert binary records
          mit Anreicherung (Profil-id, paper\_ref, axes)
    \item \texttt{csv\_to\_latex} --- generiert booktabs-Tabellen +
          Algorithmus-Steckbriefe pro Profil
    \item \texttt{diagram\_generator} --- A4-aware TikZ-Diagramme
          (Bar-Chart, Scatter, Heatmap)
\end{enumerate}

Diese Tools werden vom Manuskript-Build (\texttt{thesis/}) konsumiert
ueber \texttt{\textbackslash input\{tikz/<spec\_id>/...\}} bzw.
\texttt{\textbackslash input\{tabellen/<spec\_id>...\}}.

\section{Hypothesen}

\subsection{H1: Page-Type-Cost}

\textbf{Hypothese:} Die mittleren Zugriffskosten pro Page-Type-Variante
unterscheiden sich systematisch nach Workload-Charakteristik. Insbesondere
wird erwartet, dass dichte Page-Typen (DenseByteArt256, HotMultiByte)
unter Zipfian-Workloads (YCSB-A) schneller sind als sparse Varianten
(B+-Node), waehrend Range-Scan-optimierte Pages (B$^2$-Tree, Wormhole)
unter YCSB-E fuehren.

\textbf{Mess-Mechanik:} \texttt{H1PageTypeCost::note\_lookup(page\_type,
cycles)} in \texttt{prt\_art/measurement/hypothesis\_metrics.hpp}
sammelt online-Mean-Werte pro Page-Type ueber alle Lookup-Operationen.

\subsection{H2: Code-Qualitaet pro Bausteine-Quelle}

\textbf{Hypothese:} Der vom Repo-Originalcode geerbte Code-Qualitaets-
Score korreliert messbar mit der erreichten Throughput-Performance.
Habich-Direktive (Sprechstunde 2026-05-08) verlangt explizit die
Bewertung der Code-Qualitaet pro Bausteine-Quelle.

\textbf{Mess-Mechanik:} \texttt{H2CodeQuality::score(source, value, note)}
sammelt manuell-vergebene + automatisch-gemessene Quality-Scores pro
Source-Repo-Identifier (P01-unodb, P02-hot, A04-mimalloc, etc.).

\subsection{H3: Inline-vs-External-vs-Chain Distribution}

\textbf{Hypothese:} Die empirisch beobachtete Verteilung der drei
ValueHandle-Varianten (Inline / External / ChainRef) korreliert mit
der Page-Density: dichte Pages bevorzugen Inline (kein Indirection-
Overhead), spaerliche Pages External (Cache-Footprint-Reduktion),
PRT-ART-spezifische Heuristiken bevorzugen ChainRef (Linear-Value-
Buffer-Sharing).

\textbf{Mess-Mechanik:}
\texttt{H3InlineVsExternalDistribution::note\_inline}/\texttt{note\_external}/
\texttt{note\_chain} zaehlen pro Insert-Operation den gewaehlten
ValueHandle-Variant. Die \texttt{inline\_share}/\texttt{external\_share}/
\texttt{chain\_share}-Methoden liefern die normalisierten Anteile.

\section{Notify-Hooks (V11.1 verdrahtet)}

Die \texttt{search\_engine}-ABI definiert drei
Such-Heuristik-Hooks, die vom \texttt{ExperimentDriver} oder externen
Beobachtern aktiviert werden koennen:

\begin{itemize}
    \item \texttt{notify\_density\_threshold(pct)} $\to$
          \texttt{DensityTracker::note\_observation(pct/100)}
    \item \texttt{notify\_hot\_path\_detected(path)} $\to$
          \texttt{PathOrientedPrefetch::note\_hot\_path\_bytes(path.data,
          path.size)}
    \item \texttt{notify\_workload\_change(locality)} $\to$
          \texttt{HypothesisMetrics::adapt\_locality(locality)}
\end{itemize}

Diese Hooks sind in \texttt{PrtArtSearchEngineAdapterMap} (V11.1)
operational verdrahtet --- damit fliessen externe Beobachtungen direkt
in die H1/H2/H3-Mess-Pipeline.

\section{Mess-Resultate (TBD)}

% V15.2 / V15.3: hier kommen nach den ersten messung_driver-Laeufen die
% TikZ-Diagramme + Tabellen rein. Beispielhafte Struktur:
%
% \subsection{Messreihe A\_full}
% \input{tikz/A_full/cycles_per_permutation.tikz}
% \input{tabellen/A_full_summary.tex}
%
% \subsection{Messreihe B}
% \input{tikz/B_CacheEngine_Perms/heatmap.tikz}
% ...
